\section{The Unimodular Group}

Let n be a positive integer. We denote by \(\mathbf{SL}_{n}(\mathrm{D})\) the kernel of the group homomorphism \(\det: \mathbf{GL}_{n}(\mathrm{D}) \to \mathrm{D}_{\mathrm{ab}}^{*}\) , and we call it the unimodular group or special linear group (cf.~III, §8, No.~9, p.~537 when the field D is commutative).

THEOREM 1. --- Suppose \(n \geqslant 2\) .

\begin{enumerate}
\def\labelenumi{\alph{enumi})}
\item
  The subgroup \(\mathbf{SL}_n(\mathrm{D})\) of \(\mathbf{GL}_n(\mathrm{D})\) is generated by the matrices \(B_{ij}(\lambda)\) , where \(i\) and \(j\) are distinct integers in the interval \([1,n]\) and \(\lambda\) runs through D.
\item
  Suppose \(n \geqslant 3\) or \(\operatorname{Card}(\mathrm{D}) \geqslant 3\) . The derived group of \(\mathbf{GL}_n(\mathrm{D})\) is equal to \(\mathbf{SL}_n(\mathrm{D})\) .
\item
  Suppose \(n \geqslant 3\) or \(\operatorname{Card}(\mathrm{D}) \geqslant 4\) . The derived group of \(\mathbf{SL}_n(\mathrm{D})\) is equal to \(\mathbf{SL}_n(\mathrm{D})\) .
\end{enumerate}

\begin{enumerate}
\def\labelenumi{\Alph{enumi})}
\tightlist
\item
  Denote by T the subgroup of \(\mathbf{GL}_{n}(\mathrm{D})\) generated by the matrices \(B_{ij}(\lambda)\) . By Example 1 of VIII, p.~454, we have \(\det(B_{ij}(\lambda)) = 1\) , and therefore \(\mathrm{T} \subset \mathbf{SL}_{n}(\mathrm{D})\) . To prove that these two groups are equal, it then suffices, by Example 3 of loc. cit., to prove that every matrix of the form \(\operatorname{diag}(1, \ldots, 1, d)\) with \(\pi(d) = 1\) belongs to T. The matrix \(\operatorname{diag}(1, \ldots, 1, d)\) belongs to the image of the homomorphism \(U \mapsto \begin{pmatrix} I_{n-2} & 0 \\ 0 & U \end{pmatrix}\) from \(\mathbf{GL}_{2}(\mathrm{D})\) to \(\mathbf{GL}_{n}(\mathrm{D})\) ; it therefore suffices to consider the case n = 2. Since the kernel of \(\pi\) is the derived group of \(D^{*}\) , we may assume \(d = uvu^{-1}v^{-1}\) with u, v in \(D^{*}\) . Our assertion then follows from the equalities
\end{enumerate}

\[
\left( \begin{array}{c c} 1 & 0 \\ 0 & d \end{array} \right) = \left( \begin{array}{c c} u ^ {- 1} & 0 \\ 0 & u \end{array} \right) \left( \begin{array}{c c} v ^ {- 1} & 0 \\ 0 & v \end{array} \right) \left( \begin{array}{c c} v u & 0 \\ 0 & u ^ {- 1} v ^ {- 1} \end{array} \right)\tag{28}
\]

and

\[
\left( \begin{array}{c c} s & 0 \\ 0 & s ^ {- 1} \end{array} \right) = \left( \begin{array}{c c} 1 & s \\ 0 & 1 \end{array} \right) \left( \begin{array}{c c} 1 & 0 \\ 1 - s ^ {- 1} & 1 \end{array} \right) \left( \begin{array}{c c} 1 & - 1 \\ 0 & 1 \end{array} \right) \left( \begin{array}{c c} 1 & 0 \\ 1 - s & 1 \end{array} \right)\tag{29}
\]

for \(s \in D^{*}\) .

\begin{enumerate}
\def\labelenumi{\Alph{enumi})}
\setcounter{enumi}{1}
\tightlist
\item
  By construction, \(\mathbf{SL}_{n}(\mathrm{D})\) is the kernel of a homomorphism from \(\mathbf{GL}_{n}(\mathrm{D})\) to a commutative group, so it contains the derived group \((\mathbf{GL}_{n}(\mathrm{D}),\mathbf{GL}_{n}(\mathrm{D}))\) of \(\mathbf{GL}_{n}(\mathrm{D})\) . By the above, we have
\end{enumerate}

\[
\mathbf {S L} _ {n} (\mathrm{D}) \supset (\mathbf {G L} _ {n} (\mathrm{D}), \mathbf {G L} _ {n} (\mathrm{D})) \supset (\mathbf {S L} _ {n} (\mathrm{D}), \mathbf {S L} _ {n} (\mathrm{D})).
\]

To prove c), it suffices to prove that the matrices \(B_{ij}(\lambda)\) are commutators in \(\mathbf{SL}_n(\mathrm{D})\) .

Suppose \(n \geqslant 3\) . If \(i, j, k\) are distinct integers in the interval \([1, n]\) and \(\mu, \nu\) are elements of \(D\) , then we have

\[
B _ {i j} (\mu \nu) = B _ {i k} (\mu) ^ {- 1} B _ {k j} (\nu) ^ {- 1} B _ {i k} (\mu) B _ {k j} (\nu).\tag{30}
\]

By taking \(\mu = 1\) and \(\nu = \lambda\) , we see that the matrix \(B_{ij}(\lambda)\) is a commutator of elements of \(\mathbf{SL}_n(\mathrm{D})\) .

Now suppose \(n = 2\) . Let \(u\) and \(v\) be elements of \(D\) with \(u \neq 0\) . We have the relations

\[
\left( \begin{array}{c c} 1 & v - u v u \\ 0 & 1 \end{array} \right) = \left( \begin{array}{c c} u & 0 \\ 0 & u ^ {- 1} \end{array} \right) \left( \begin{array}{c c} 1 & - v \\ 0 & 1 \end{array} \right) \left( \begin{array}{c c} u ^ {- 1} & 0 \\ 0 & u \end{array} \right) \left( \begin{array}{c c} 1 & v \\ 0 & 1 \end{array} \right)\tag{31}
\]

and

\[
\left( \begin{array}{c c} 1 & 0 \\ v - u v u & 1 \end{array} \right) = \left( \begin{array}{c c} u ^ {- 1} & 0 \\ 0 & u \end{array} \right) \left( \begin{array}{c c} 1 & 0 \\ - v & 1 \end{array} \right) \left( \begin{array}{c c} u & 0 \\ 0 & u ^ {- 1} \end{array} \right) \left( \begin{array}{c c} 1 & 0 \\ v & 1 \end{array} \right).\tag{32}
\]

We have \(\det\left(\begin{matrix}u & 0 \\ 0 & u^{-1}\end{matrix}\right)=1\) , so the matrices \(B_{12}(v-uvu)\) and \(B_{21}(v-uvu)\) are commutators of elements of \(\mathbf{SL}_{n}(\mathrm{D})\) .

Suppose that the field D has at least four elements. Let \(\lambda\) be an element of D. If \(\lambda\) is equal to 0, 1, or -1, choose an arbitrary element u in \(D-\{0,1,-1\}\) ; otherwise, set \(u=\lambda\) . In both cases, u is a nonzero element of D, it commutes with \(\lambda\) , and we have \(u^{2}\neq1\) . Set \(v=\lambda(1-u^{2})^{-1}\) . We have uv=vu and therefore \(v-uvu=v(1-u^{2})=\lambda\) . It then follows from relations (31) and (32) that the matrices \(B_{12}(\lambda)\) and \(B_{21}(\lambda)\) are commutators in \(\mathbf{SL}_{n}(\mathrm{D})\) . This completes the proof of c).

\begin{enumerate}
\def\labelenumi{\Alph{enumi})}
\setcounter{enumi}{2}
\tightlist
\item
  It remains to prove that \(\mathbf{SL}_{2}(\mathrm{D})\) is the derived group of \(\mathbf{GL}_{2}(\mathrm{D})\) when D has three elements. By A), the group \(\mathbf{SL}_{2}(\mathrm{D})\) is generated by the matrices \(B_{12}(1)=(\begin{matrix}{1}&{1}\\ {0}&{1}\end{matrix})\) and \(B_{21}(1)=(\begin{matrix}{1}&{0}\\ {1}&{1}\end{matrix})\) , and these matrices are commutators of elements of \(\mathbf{GL}_{2}(\mathrm{D})\) because we have \(B_{21}(1)={}^{t}B_{12}(1)\) and
\end{enumerate}

\[
\left( \begin{array}{c c} 1 & 1 \\ 0 & 1 \end{array} \right) = \left( \begin{array}{c c} - 1 & 0 \\ 0 & 1 \end{array} \right) \left( \begin{array}{c c} 1 & 1 \\ 0 & 1 \end{array} \right) \left( \begin{array}{c c} - 1 & 0 \\ 0 & 1 \end{array} \right) ^ {- 1} \left( \begin{array}{c c} 1 & 1 \\ 0 & 1 \end{array} \right) ^ {- 1}.\tag{33}
\]

Remark. --- If D is a field with two elements, then \(\mathbf{GL}_{2}(\mathrm{D})\) is equal to \(\mathbf{SL}_{2}(\mathrm{D})\) , and it is a group of order 6 whose derived group has order 3. If D is a field with three elements, then the group \(\mathbf{SL}_{2}(\mathrm{D})\) has order 24, and its derived group has order 8. Suppose \(n \geqslant 2\) ; except in the two previous cases, every normal subgroup of \(\mathbf{SL}_{n}(\mathrm{D})\) distinct from \(\mathbf{SL}_{n}(\mathrm{D})\) is contained in the center \(\mathrm{Z}(\mathbf{SL}_{n}(\mathrm{D}))\) of \(\mathbf{SL}_{n}(\mathrm{D})\) (II, §10, p.~421, Exercises 13 and 14). The group \(\mathbf{SL}_{n}(\mathrm{D})/\mathrm{Z}(\mathbf{SL}_{n}(\mathrm{D}))\) is then simple.

Let V be a right vector space over the field D, of finite dimension n.~We denote by \(\mathrm{SL}(V)\) , and call unimodular group of V, the kernel of the homomorphism \(\det: \mathbf{GL}(V) \to D_{ab}^{*}\) . Choosing a basis of V allows one to identify this group with \(\mathbf{SL}_{n}(D)\) .

We say that an automorphism u of V is a transvection if there exist a vector v of V and a linear form \(\varphi\) on V such that \(\varphi(v)=0\) and \(u(x)=x+v\varphi(x)\) for every x in V. When \(n\geqslant2\) , u is a transvection if and only if there exists a basis of V with respect to which the matrix of u is of the form \(B_{ij}(\lambda)\) . Theorem 1 implies the following corollary.

COROLLARY. --- Let V be a right vector space over the field D, of finite dimension \(n \geqslant 2\) .

\begin{enumerate}
\def\labelenumi{\alph{enumi})}
\item
  The subgroup \(\mathbf{SL}(\mathrm{V})\) of \(\mathbf{GL}(\mathrm{V})\) is generated by the transvections.
\item
  The subgroup \(\mathbf{SL}(\mathrm{V})\) is the derived group of \(\mathbf{GL}(\mathrm{V})\) unless we have \(n = 2\) and D has two elements.
\item
  The group \(\mathbf{SL}(V)\) is equal to its derived group unless we have n = 2 and D has two or three elements.
\end{enumerate}

